
\documentclass{article}
\usepackage{colortbl}
\usepackage{makecell}
\usepackage{multirow}
\usepackage{supertabular}

\begin{document}

\newcounter{utterance}

\centering \large Interaction Transcript for game `wordle\_withcritic', experiment `medium\_frequency\_words\_clue\_with\_critic', episode 3 with reflexion{-}llama{-}3.1{-}8b{-}instruct.
\vspace{24pt}

{ \footnotesize  \setcounter{utterance}{1}
\setlength{\tabcolsep}{0pt}
\begin{supertabular}{c@{$\;$}|p{.15\linewidth}@{}p{.15\linewidth}p{.15\linewidth}p{.15\linewidth}p{.15\linewidth}p{.15\linewidth}}
    \# & $\;$A & \multicolumn{4}{c}{Game Master} & $\;\:$B\\
    \hline

    \theutterance \stepcounter{utterance}  
    & & \multicolumn{4}{p{0.6\linewidth}}{
        \cellcolor[rgb]{0.9,0.9,0.9}{
            \makecell[{{p{\linewidth}}}]{
                \texttt{\tiny{[P1$\langle$GM]}}
                \texttt{You have attempted this same task before and did not succeed. These are the plans you wrote after those attempts.} \\
\texttt{Trial \#1: New plan: My last attempt got aborted because I produced a guess that was too short (only 4 letters). I should have paid closer attention to the length of the word "sweetheart" and taken that into account when making my guess. Next time, I will make sure to produce a 5{-}letter word that is related to the clue "sweetheart" and the theme of affection or love. For example, I could try a word like "honey".} \\
\\ 
\texttt{Trial \#2: New plan: My last attempt got aborted because I produced a guess that was too short (only 4 letters) again, despite my previous plan to pay closer attention to the length of the word "sweetheart". I should have noticed that the Game Master agreed with my previous guess "honey" which is a 5{-}letter word, and I should have built upon that. Next time, I will make sure to produce a 5{-}letter word that is related to the clue "sweetheart" and the theme of affection or love, and I will also pay attention to the length of the word. For example, I could try a word like "lovely".} \\
\\ 
\\ 
\texttt{You are a language wizard who likes to guess words by using the given rules.} \\
\\ 
\texttt{Welcome to Wordle! You have six attempts to guess the target word, a valid English word of five lowercase letters (a{-}z). Please use the tags "explanation:" and "guess:" to provide a concise explanation for each guess.} \\
\\ 
\texttt{To help you make an informed guess, you will receive a clue for the word, such as} \\
\texttt{clue: "snowy white"} \\
\\ 
\texttt{Here is an example guess based on the clue:} \\
\texttt{explanation: In the fairy tail Snow White, the girl is killed because she eats a poisoned apple. And the word apple has 5 letters.} \\
\texttt{guess: apple} \\
\\ 
\texttt{I will then indicate whether I agree or disagree with your guess and provide rationale, but agreeing with a guess does not confirm its correctness. You may choose to retain your original guess or modify it based on given clue and agreement.} \\
\\ 
\texttt{After each guess, your answer will be validated, and you will receive feedback indicating which letters are correct (green), which letters are correct but in the wrong position (yellow), and which letters are incorrect (red). This feedback can be useful in determining which letters to include or exclude in your next guess.} \\
\\ 
\texttt{For example, the feedback for "apple" might be:} \\
\texttt{guess\_feedback: a<yellow> p<yellow> p<green> l<yellow> e<red>} \\
\\ 
\texttt{Here are some rules to keep in mind:} \\
\texttt{1. Please do not generate any guess feedback or clue on your own} \\
\texttt{2. Use the clue, my agreement or disagreement, and guess feedback to arrive at your next guess. This should be reflected in your explanation as well.} \\
\\ 
\texttt{Do not generate any other text.} \\
\\ 
\texttt{Let's begin with your first guess.} \\
\\ 
\texttt{clue: sweetheart} \\
            }
        }
    }
    & \\ \\

    \theutterance \stepcounter{utterance}  
    & \multicolumn{4}{p{0.6\linewidth}}{
        \cellcolor[rgb]{0.9,0.9,0.9}{
            \makecell[{{p{\linewidth}}}]{
                \texttt{\tiny{[P1$\rangle$GM]}}
                \texttt{explanation: The clue is "sweetheart" which is a term of endearment. I will try a 5{-}letter word related to affection or love. } \\
\texttt{guess: honey} \\
            }
        }
    }
    & & \\ \\

    \theutterance \stepcounter{utterance}  
    & & \multicolumn{4}{p{0.6\linewidth}}{
        \cellcolor[rgb]{0.9,0.9,0.9}{
            \makecell[{{p{\linewidth}}}]{
                \texttt{\tiny{[GM$\rangle$P2]}}
                \texttt{You have attempted this same task before and did not succeed. These are the plans you wrote after those attempts.} \\
\texttt{Trial \#1: New plan: I got the agreement right, but the format was wrong. I should have stuck to the format and only provided "yes" or "no" as the agreement, without any additional text. Next time, I will make sure to follow the format exactly and only provide the required information.} \\
\\ 
\\ 
\texttt{I need your assistance with a word game in which we need to find a 5{-}letter word using a clue, a guess and an explanation for the guess.} \\
\\ 
\texttt{For example, if the clue is "snow white" and the guess is "apple", you will get the information like this:} \\
\texttt{clue: snow white} \\
\texttt{explanation: the fairy tale Snow White involves the girl eating a poisoned apple, which has 5 letters.} \\
\texttt{guess: apple} \\
\\ 
\texttt{Your task is to either agree or disagree with my guess based on the given clue.} \\
\\ 
\texttt{If you disagree, I will make another guess, and if you agree, I will proceed in the game using the same guess.} \\
\\ 
\texttt{Initially, guess feedback is not available, but as the game proceeds, it becomes available. The guess feedback is not for current guess, but rather an indication of what letters can be present in the current guess. A letter highlighted in green indicates that it is a correct letter in the correct position, while yellow indicates that it is a correct letter in the wrong position, and red indicates a wrong letter. At this point, you should use all the provided information, including the clue, guess, explanation, and guess feedback, to decide whether the given guess matches the clue, and also stick to guess feedback or not.} \\
\\ 
\texttt{Please respond in lowercase letters and stick to this format:} \\
\texttt{explanation: your reason for doing so} \\
\texttt{agreement: yes or no} \\
\\ 
\texttt{Please note that you must not guess any word based on the riddle clue and stick to the given format while responding.} \\
\\ 
\texttt{Do not generate any other text.} \\
\\ 
\texttt{Let's begin!} \\
\\ 
\texttt{clue: sweetheart} \\
\texttt{explanation: The clue is "sweetheart" which is a term of endearment. I will try a 5{-}letter word related to affection or love.} \\
\texttt{guess: honey} \\
            }
        }
    }
    & \\ \\

    \theutterance \stepcounter{utterance}  
    & & & \multicolumn{4}{p{0.6\linewidth}}{
        \cellcolor[rgb]{0.9,0.9,0.9}{
            \makecell[{{p{\linewidth}}}]{
                \texttt{\tiny{[GM$\langle$P2]}}
                \texttt{explanation: the guess "honey" matches the clue "sweetheart" as honey is a term of endearment and is related to affection or love.} \\
\texttt{agreement: yes} \\
            }
        }
    }
     \\ \\

    \theutterance \stepcounter{utterance}  
    & & \multicolumn{4}{p{0.6\linewidth}}{
        \cellcolor[rgb]{0.9,0.9,0.9}{
            \makecell[{{p{\linewidth}}}]{
                \texttt{\tiny{[P1$\langle$GM]}}
                \texttt{clue: sweetheart} \\
\texttt{guess\_agreement: yes} \\
\texttt{agreement\_explanation: the guess "honey" matches the clue "sweetheart" as honey is a term of endearment and is related to affection or love.} \\
\\ 
\texttt{Provide your response only in this format.} \\
\texttt{explanation: details} \\
\texttt{guess: word} \\
            }
        }
    }
    & \\ \\

    \theutterance \stepcounter{utterance}  
    & \multicolumn{4}{p{0.6\linewidth}}{
        \cellcolor[rgb]{0.9,0.9,0.9}{
            \makecell[{{p{\linewidth}}}]{
                \texttt{\tiny{[P1$\rangle$GM]}}
                \texttt{explanation: Since the Game Master agreed with my previous guess "honey" which is a 5{-}letter word, I will try another 5{-}letter word related to affection or love.} \\
\texttt{guess: lovely} \\
            }
        }
    }
    & & \\ \\

    \theutterance \stepcounter{utterance}  
    & & & \multicolumn{2}{p{0.3\linewidth}}{
        \cellcolor[rgb]{0.9,0.9,0.9}{
            \makecell[{{p{\linewidth}}}]{
                \texttt{\tiny{[GM$|$GM]}}
                \texttt{Error: The length of the guessed word is not 5.} \\
            }
        }
    }
    & & \\ \\

    \theutterance \stepcounter{utterance}  
    & & & \multicolumn{2}{p{0.3\linewidth}}{
        \cellcolor[rgb]{0.9,0.9,0.9}{
            \makecell[{{p{\linewidth}}}]{
                \texttt{\tiny{[GM$|$GM]}}
                \texttt{game\_result = ABORT} \\
            }
        }
    }
    & & \\ \\

    \theutterance \stepcounter{utterance}  
    & \multicolumn{4}{p{0.6\linewidth}}{
        \cellcolor[rgb]{0.95,0.95,0.95}{
            \makecell[{{p{\linewidth}}}]{
                \texttt{\tiny{[P2$\rightarrow$P1]}}
                \texttt{Post{-}episode Reflexion (not sent to GM):} \\
\texttt{New plan: My last two attempts got aborted because I produced a guess that was too short (only 4 letters) despite my previous plans to pay closer attention to the length of the word "sweetheart". I should have noticed that the Game Master agreed with my previous guess "honey" which is a 5{-}letter word, and I should have built upon that. Next time, I will make sure to produce a 5{-}letter word that is related to the clue "sweetheart" and the theme of affection or love, and I will also pay attention to the length of the word. However, I should also pay closer attention to the Game Master's agreement and use it to inform my next guess. For example, if the Game Master agrees with a 5{-}letter word, I should try to modify it or build upon it instead of starting from scratch.} \\
            }
        }
    }
    & & \\ \\

    \theutterance \stepcounter{utterance}  
    & & & \multicolumn{4}{p{0.6\linewidth}}{
        \cellcolor[rgb]{0.95,0.95,0.95}{
            \makecell[{{p{\linewidth}}}]{
                \texttt{\tiny{[P1$\rightarrow$P2]}}
                \texttt{Post{-}episode Reflexion (not sent to GM):} \\
\texttt{New plan: I got the agreement right, but the format was wrong. I should have stuck to the format and only provided "yes" or "no" as the agreement, without any additional text. Next time, I will make sure to follow the format exactly and only provide the required information. I should have responded with "agreement: yes" only, without the explanation.} \\
            }
        }
    }
     \\ \\

\end{supertabular}
}

\end{document}
